\section*{Chapitre \ref{ch:g9:natural-risks} --- Vivre avec les risques naturels}

\begin{solution}{exo:g9:natural-risks:1}
(a) Aléa. (b) \omterm{def:g9:natural-risks:vulnerability}{Vulnérabilité}. (c) \omterm{def:g9:natural-risks:vulnerability}{Vulnérabilité}. (d) Aléa. (e) Les deux~:
les coulées de boue sont l'aléa~; bâtir la ville sur leur trajet la rend
vulnérable.
\end{solution}

\begin{solution}{exo:g9:natural-risks:2}
(a) Faux~: là où personne ne vit, rien ne peut subir de dommages, donc le
\omterm{def:g9:natural-risks:risk}{risque} est faible. (b) Vrai. (c) Vrai. (d) Faux~: une alerte ne change
pas le \omterm{def:g4:volcanoes-earthquakes-first-look:volcano}{volcan}~; elle réduit la \omterm{def:g9:natural-risks:vulnerability}{vulnérabilité} des habitants, qui peuvent
partir à temps.
\end{solution}

\begin{solution}{exo:g9:natural-risks:3}
(a) Oui, quelques jours à l'avance. (b) Oui, quelques heures à quelques
jours à l'avance. (c) Non~: les \omterm{def:g4:volcanoes-earthquakes-first-look:earthquake}{séismes} ne peuvent pas être prévus.
(d) Oui~: une fois le \omterm{def:g7:earthquakes:tsunami}{tsunami} déclenché, il met des heures à traverser
l'\omterm{def:g2:water-on-earth:ocean}{océan}, et l'alerte voyage plus vite.
\end{solution}

\begin{solution}{exo:g9:natural-risks:4}
(a) Les secousses du sol~: des essaims de petits \omterm{def:g4:volcanoes-earthquakes-first-look:earthquake}{séismes} quand le \omterm{def:g7:volcanoes:magma}{magma}
monte. (b) L'inclinaison de la pente~: le gonflement du \omterm{def:g4:volcanoes-earthquakes-first-look:volcano}{volcan}. (c) Le
déplacement du sol, de quelques millimètres~: le gonflement quand le
\omterm{def:g7:volcanoes:magma}{magma} s'accumule. (d) Les gaz du panache~: davantage de gaz s'échappe
quand le \omterm{def:g7:volcanoes:magma}{magma} monte.
\end{solution}

\begin{solution}{exo:g9:natural-risks:5}
(a) $45 \div 2 = \qty{22.5}{km/h}$. (b) $4 \times 2 = \qty{8}{km}$, assez
pour atteindre des terrains plus élevés. (c) Personne n'a vu l'\omterm{def:g4:volcanoes-earthquakes-first-look:eruption}{éruption},
cachée par l'orage, et personne ne les a prévenus~: il n'existait aucun
moyen rapide de transmettre l'information à la ville, ni aucun plan
d'évacuation. (d) Par exemple~: un plan d'évacuation avec des itinéraires
vers les hauteurs, répété à l'avance~; une chaîne d'alerte allant des
scientifiques à la ville~; davantage d'instruments sur le \omterm{def:g4:volcanoes-earthquakes-first-look:volcano}{volcan}.
\end{solution}

\begin{solution}{exo:g9:natural-risks:6}
(a) La \omterm{def:g5:water-shapes-the-land:flood}{crue} de \qty{10}{m} au-dessus de la \omterm{def:g1:landscapes-around-us:water}{rivière}. (b) Environ
\qty{0.47}{km} pour la \omterm{def:g5:water-shapes-the-land:flood}{crue} de \qty{2}{m}, \qty{0.56}{km} pour la \omterm{def:g5:water-shapes-the-land:flood}{crue} de
\qty{5}{m} et \qty{0.67}{km} pour la \omterm{def:g5:water-shapes-the-land:flood}{crue} de \qty{10}{m} (une lecture à
environ \qty{0.05}{km} près convient). (c) Seulement la \omterm{def:g5:water-shapes-the-land:flood}{crue} de
\qty{10}{m}~: le terrain est plus haut que 2 et \qty{5}{m} au-dessus de
la \omterm{def:g1:landscapes-around-us:water}{rivière}. (d) Il suffit d'une petite montée de la \omterm{def:g1:landscapes-around-us:water}{rivière} pour
l'atteindre, et les petites \omterm{def:g5:water-shapes-the-land:flood}{crues} reviennent plus souvent que les
grandes.
\end{solution}

\begin{solution}{exo:g9:natural-risks:7}
(a) Aléa fort, \omterm{def:g9:natural-risks:vulnerability}{vulnérabilité} faible~: \omterm{def:g9:natural-risks:risk}{risque} moyen. (b) Aléa faible,
\omterm{def:g9:natural-risks:vulnerability}{vulnérabilité} forte~: \omterm{def:g9:natural-risks:risk}{risque} moyen. (c) Aléa fort, \omterm{def:g9:natural-risks:vulnerability}{vulnérabilité} forte~:
\omterm{def:g9:natural-risks:risk}{risque} élevé. (d) Aléa faible, \omterm{def:g9:natural-risks:vulnerability}{vulnérabilité} faible~: \omterm{def:g9:natural-risks:risk}{risque} faible.
\end{solution}

\begin{solution}{exo:g9:natural-risks:8}
(a) C, puis B, puis A. (b) Les contreventements forment des triangles,
qui ne peuvent pas se déformer sans étirer ou écraser une barre~:
l'ossature garde sa forme au lieu de pencher. (c) Les isolateurs cèdent
quand le sol va et vient sous eux, si bien qu'une bien plus petite part
du mouvement atteint l'école.
\end{solution}

\begin{solution}{exo:g9:natural-risks:9}
(a) $1 \div 100 = 1~\%$ par an. (b) La chance est la même chaque année,
quoi qu'il se soit passé l'année précédente~: la \omterm{def:g5:water-shapes-the-land:flood}{crue} peut revenir dès
l'an prochain. (c) $300 \div 100 = 3$ \omterm{def:g5:water-shapes-the-land:flood}{crues}, en moyenne.
\end{solution}

\begin{solution}{exo:g9:natural-risks:10}
(a) $3600 \div 15 = 240$ mesures. (b) $240 \times 24 \times 365 =
\num{2102400} \approx 2.1 \times 10^6$ mesures. (c) $2400 \div 800 =
\qty{3}{h}$. (d) Pendant ces trois heures, les centres d'alerte peuvent
prévenir la \omterm{def:g1:landscapes-around-us:water}{côte} et la population peut gagner les hauteurs.
\end{solution}

\begin{solution}{exo:g9:natural-risks:11}
(a) $4 \times 3 \times 3.8 = \qty{45.6}{L}$. (b) $45.6 \div 1.5 = 30.4$,
donc 31 bouteilles. (c) Par exemple~: de la nourriture qui se conserve,
pour manger quand les magasins sont fermés~; une lampe de poche, si
l'électricité est coupée~; une radio à piles ou à manivelle, pour
entendre les alertes~; une trousse de premiers secours, pour soigner les
petites blessures~; des couvertures chaudes. (d) Pour pouvoir l'emporter
tout de suite si la famille doit partir vite.
\end{solution}

\begin{solution}{exo:g9:natural-risks:12}
Par exemple~: (1) établir une \omterm{def:g9:natural-risks:hazard-map}{carte des aléas} des coulées de boue, \omterm{def:g1:landscapes-around-us:landform}{vallée}
par \omterm{def:g1:landscapes-around-us:landform}{vallée} (côté aléa)~; (2) installer sur le \omterm{def:g4:volcanoes-earthquakes-first-look:volcano}{volcan} des \omterm{def:g7:earthquakes:seismometer}{sismomètres}, des
inclinomètres, des récepteurs GPS et des analyseurs de gaz (côté aléa)~;
(3) mettre en place une chaîne d'alerte rapide entre l'observatoire et la
ville, avec des sirènes (côté \omterm{def:g9:natural-risks:vulnerability}{vulnérabilité})~; (4) prévoir et baliser des
itinéraires vers les hauteurs, et s'y entraîner par des exercices (côté
\omterm{def:g9:natural-risks:vulnerability}{vulnérabilité})~; (5) ne construire aucune nouvelle habitation ni école sur
le trajet des coulées de boue (côté \omterm{def:g9:natural-risks:vulnerability}{vulnérabilité}).
\end{solution}

\begin{solution}{pb:g9:natural-risks:1}
\textbf{1.} La \omterm{def:g1:landscapes-around-us:water}{rivière} est à \qty{143}{m} à côté de C, à \qty{137}{m} à
côté de A et à \qty{126}{m} à côté de B~; sur la \omterm{def:g2:drawing-to-map:map}{carte}, C est à l'ouest et
B à l'est, donc la \omterm{def:g1:landscapes-around-us:water}{rivière} coule vers l'est.

\textbf{2.} A~: $137 - 137 = \qty{0}{m}$~; B~: $276 - 126 = \qty{150}{m}$~;
C~: $236 - 143 = \qty{93}{m}$.

\textbf{3.} Seulement A, qui est à moins de \qty{5}{m} au-dessus de la \omterm{def:g1:landscapes-around-us:water}{rivière}.

\textbf{4.} Environ \qty{0.57}{km} (entre 0.5 et \qty{0.65}{km}).

\textbf{5.} L'élévation sur \qty{100}{m} est la moitié de l'élévation sur
\qty{200}{m}~: A $1.4 \div 2 = 0.7~\%$~; B $83 \div 2 = 41.5~\%$~;
C $20 \div 2 = 10~\%$.

\textbf{6.} B~: sa pente est très raide, et les \omterm{def:g3:stones-pebbles-sand:rock}{roches} et la terre
meubles peuvent glisser le long d'une pente raide, surtout après de
fortes \omterm{def:g4:clouds-rain-snow:precipitation}{pluies}.

\textbf{7.} A~: environ \qty{1}{km}~; B~: presque sur la \omterm{def:g7:earthquakes:fault}{faille} (moins de
\qty{0.05}{km})~; C~: environ \qty{1.3}{km}.

\textbf{8.} Si la \omterm{def:g7:earthquakes:fault}{faille} joue, le sol se rompt et se décale le long de
celle-ci, ce qui peut déchirer un bâtiment~; et un \omterm{def:g4:volcanoes-earthquakes-first-look:earthquake}{séisme} sur la \omterm{def:g7:earthquakes:fault}{faille}
secouerait très fort le sol à cet endroit.

\textbf{9.}

\begin{center}
{\small
\begin{tabular}{|c|c|c|c|}
\hline
terrain & dans la zone inondable~? & pente raide~? & sur la \omterm{def:g7:earthquakes:fault}{faille}~? \\
\hline
A & oui & non & non \\
B & non & oui & oui \\
C & non & non & non \\
\hline
\end{tabular}}
\end{center}

\textbf{10.} Le terrain C~: c'est le seul qui a trois «~non~»~:
\qty{93}{m} au-dessus de la \omterm{def:g1:landscapes-around-us:water}{rivière}, une pente douce de \qty{10}{\%}, à
plus de \qty{1}{km} de la \omterm{def:g7:earthquakes:fault}{faille}.

\textbf{11.} Par exemple~: des contreventements en diagonale et des voiles
dans l'ossature, des assemblages solides, ou des isolateurs sous le
bâtiment~; les élèves s'entraînent aux exercices de \omterm{def:g4:volcanoes-earthquakes-first-look:earthquake}{séisme} (se baisser,
s'abriter, s'agripper).

\textbf{12.} \textbf{17 maisons} sur 28 se trouvent dans la zone
inondable~: $17 \div 28 \approx 0.61$, environ 3 maisons sur 5.
\end{solution}
